% Pursuing Stacks (cotes 134-2 to 134-8): archive cote 134-5, site page 117 (PDF page 118), typescript page 304.
% typescript. Draft compared: ps5.tex:615-683.
% AI-transcribed by Claude Opus 5.5 (claude-opus-5-5, adaptive thinking, effort high), Message Batches API, batch ps3 (msgbatch_01XdFmuVVBAcgc7EFxdtCqsj), 2026-11-12; not yet checked by eye.
\scanpage{134-5/117}
\mspage{304}

more stringently still\xout{\unsure{s}} of \hand{``}c-\uline{aspheric} \uline{functors}\hand{''}(as we may call them\over{\unsure{l}}{)} is
rather evident,  as it gives a way,  via $i^{*}$ ,  for any homotopy type
described in terms of a ``model'' x in M,  to find a corresponding model
$i^{*}(x)$ in $A^{\wedge}$ .  The situation would be more satisfactory still if we \xout{\unsure{coul}}
could find the test functor i such that the corresponding functor
\[ \hand{(1)} \qquad i_{!} : A^{\wedge} \longrightarrow M \]
be modelizing too  (assuming M to be stable under small direct limits,
so that $i_{!}$ is defined as the left adjoint of \xout{\illegible{1}}
\[ \hand{(2)} \qquad i^{*} : M \longrightarrow A^{\wedge} \quad .) \]
In this case,  for a homotopy type described by a model \deleted{\unsure{X}}\inserted{K} in $A^{\wedge}$,
$i_{!}(K)$  gives a description of the same by a model in M.

Maybe we should remember though that even if we do not know about
any test functor from A to M,  still we always can find in three steps
a modelizing functor
\[ M \longrightarrow A^{\wedge} \ , \]
namely a composition
\[ (*) \qquad M \xrightarrow{\hand{$j$}^{*}} B^{\wedge} \xrightarrow{\hand{$i$}_{B}} (\mathrm{Cat}) \xrightarrow{j_{A} = i_{A}^{*}} A^{\wedge} \quad , \]
where  $\hand{$j$}: B \to M$  is an $M_{a}$-aspheric functor from an auxiliary small
category,  which  we may assume to be a \xout{weak} test category\hand{,} \xout{\unsure{itexxsuch}}
\deleted{as}\inserted{by} a mild extra assumption on \over{\unsure{M}}{M} (cf cor. 3  p. 253).  The modelizing
functor we thus get has  the disadvantage of not being left exact,
whereas the looked-for functor $i^{*}$ commutes to small inverse limits.
Still,  the composition (*)  is pretty near to being left exact,  it
commutes to fibered products (because $i_{B}$  does)  which is the next best
-  we can view it as a\deleted{n} \inserted{left} exact functor from  M  to  $A^{\wedge}_{/E}$ ,  where \unsure{E|B} is
the image in $A^{\wedge}$ of the final object of M (assuming \xout{\unsure{therexis}} $e_{M}$ exists).

There  is another advantage still of having a test functor
$i : A \longrightarrow M$ , \deleted{with respect to} \inserted{rather than} merely using (*),  namely it al\xout{b}lows us to
``enrich''  the category structure of M,  in such a\xout{x} way as to get
``external Hom's''  of objects of M,  with \hand{``}values in $A^{\wedge}$\hand{''},  by defining,
for  x,y  in M ,  the object  \uline{Hom}(A)(x,y)  of $A^{\wedge}$ as
\[ \hand{(3)} \qquad \underline{\mathrm{Hom}}(A)(x,y) \ = \ \hand{$\{$} a \longmapsto \mathrm{Hom}_{M}(i(a) \times x \, , \ y) \hand{$\}$} \ . \]
\xout{This} Such enr\xout{a}iched structure,  when  $A = \hand{$\Delta$}$\hand{\unsure{,}}  plays an important part
in the second part of Quillen's treatment of homotopical algebra,  un-%
der the name of \hand{(}semi-\hand{)}simplicial categories,  especially wi th the notion
of \hand{(}semi-\hand{)}simplicial model categories,  which looks quite handy indeed.

\margin{117}
